\documentclass[article]{jss}

%% -- LaTeX packages and custom commands ---------------------------------------
\usepackage{orcidlink,thumbpdf,lmodern}
\usepackage{booktabs}
\usepackage{amssymb}
\usepackage{amsmath}

\newcommand{\class}[1]{`\code{#1}'}
\newcommand{\fct}[1]{\code{#1()}}


%% -- Article metainformation -------------------------------------------------

\author{Francisco Parra\\Universidad de Santiago de Chile}
\Plainauthor{Francisco Parra}

\title{\pkg{smelt-ml}: A Pure-\proglang{Rust} Machine Learning Framework with Cache-Optimized Gradient Boosting Competitive with \proglang{C++} Implementations}
\Plaintitle{smelt-ml: A Pure-Rust Machine Learning Framework with Cache-Optimized Gradient Boosting Competitive with C++ Implementations}
\Shorttitle{\pkg{smelt-ml}: Cache-Optimized ML in \proglang{Rust}}

\Abstract{
We present \pkg{smelt-ml}, an open-source machine learning framework
implemented entirely in \proglang{Rust} that provides 27 supervised learning
algorithms, unsupervised clustering, causal inference, survival analysis,
and spatial machine learning capabilities. The framework's gradient boosting
implementations (\pkg{XGBoost}, \pkg{LightGBM}, \pkg{CatBoost}) employ a
cache-optimized column-major histogram accumulation strategy that achieves
performance competitive with --- and in many cases exceeding --- the official
\proglang{C++} implementations. On standardized benchmarks with 20 features
and up to 10{,}000 samples, our \pkg{XGBoost} is 1.1--2.6$\times$ faster,
\pkg{LightGBM} 1.2--1.4$\times$ faster, and \pkg{CatBoost} 1.3--4.2$\times$
faster than their respective \proglang{C++} counterparts. The framework is
published on \url{https://crates.io/} as \pkg{smelt-ml} and provides a
composable, trait-based architecture inspired by \pkg{mlr3}
\citep{Lang+Binder+Richter:2019} and \pkg{scikit-learn}
\citep{Pedregosa+Varoquaux+Gramfort:2011}.
}

\Keywords{machine learning, gradient boosting, \proglang{Rust}, cache optimization, \pkg{XGBoost}, \pkg{LightGBM}, \pkg{CatBoost}, software}
\Plainkeywords{machine learning, gradient boosting, Rust, cache optimization, XGBoost, LightGBM, CatBoost, software}

\Address{
  Francisco Parra\\
  Departamento de Ingenier\'ia Geogr\'afica\\
  Universidad de Santiago de Chile\\
  Avenida Libertador Bernardo O'Higgins 3363\\
  Santiago, Chile\\
  E-mail: \email{francisco.parra.o@usach.cl}
}

\begin{document}


%% -- Introduction -------------------------------------------------------------

\section{Introduction} \label{sec:intro}

Gradient boosted decision trees (GBDTs) are the dominant method for supervised
learning on tabular data \citep{Grinsztajn+Oyallon+Varoquaux:2022}. The three
most widely used implementations --- \pkg{XGBoost} \citep{Chen+Guestrin:2016},
\pkg{LightGBM} \citep{Ke+Meng+Finley:2017}, and \pkg{CatBoost}
\citep{Prokhorenkova+Gusev+Vorobev:2018} --- are implemented in \proglang{C++}
and have been optimized over many years. These libraries serve as the de facto
baselines in the majority of applied machine learning studies.

\proglang{Rust} \citep{Rust:2024} is a systems programming language that
guarantees memory safety without garbage collection through its ownership type
system \citep{Matsakis+Klock:2014}. It generates machine code comparable to
\proglang{C++} via the LLVM backend while preventing buffer overflows,
use-after-free, and data races at compile time. Despite these advantages,
\proglang{Rust}'s ML ecosystem remains nascent: \pkg{linfa} \citep{linfa:2024},
the primary \proglang{Rust} ML crate, implements only 9 algorithms and does not
include gradient boosting.

In this paper, we present \pkg{smelt-ml}, a comprehensive ML framework in
\proglang{Rust} that fills this gap. Our main contributions are:

\begin{enumerate}
\item A \textbf{complete ML framework} with 27 supervised learners, clustering,
  survival analysis, causal inference, and spatial ML --- the most comprehensive
  ML framework in \proglang{Rust}.
\item \textbf{Cache-optimized gradient boosting} using column-major bin storage
  with \code{u8} packing, achieving performance that matches or exceeds the
  official \proglang{C++} implementations.
\item \textbf{Novel algorithms} not available in competing frameworks:
  Geographical-XGBoost \citep{Grekousis:2025}, Causal Forest
  \citep{Wager+Athey:2018}, Conformal Prediction \citep{Romano+Patterson+Candes:2019},
  Oblique Forest \citep{Tomita+Browne+Shen:2020}, and streaming Hoeffding Trees.
\item \textbf{Validation on real datasets} showing accuracy comparable to
  \pkg{scikit-learn} across standard benchmarks.
\end{enumerate}


%% -- Software Architecture ---------------------------------------------------

\section{Software architecture} \label{sec:architecture}

\subsection{Design principles}

\pkg{smelt-ml} follows a trait-based architecture inspired by \pkg{mlr3}:

\begin{Code}
Data (CSV) -> Task -> Pipeline(Transformers -> Learner)
                -> Prediction -> Measure
\end{Code}

The core abstractions are defined as \proglang{Rust} traits:

\begin{itemize}
\item \code{Learner}: Algorithm that trains on a \code{Task} and produces a
  \code{TrainedModel}. Requires \code{Send + Sync} for thread safety.
\item \code{TrainedModel}: Fitted model that predicts on new data.
\item \code{Measure}: Evaluation metric (Accuracy, F1, RMSE, R$^2$, AUC-ROC, etc.).
\item \code{Transformer}: Preprocessing step composable via \code{Pipeline}.
\item \code{Resample}: Data splitting strategy.
\end{itemize}

This design enables zero-friction composition: \code{Pipeline} implements
\code{Learner}, so a pipeline of transformers and a learner can be passed to
\fct{benchmark}, \code{GridSearch}, or \code{Bagging} without adaptation.

\subsection{Quick start}

A minimal classification example:

\begin{CodeChunk}
\begin{CodeInput}
rust> use smelt_ml::prelude::*;
rust> use ndarray::array;
rust>
rust> let features = array![[0.0, 0.0], [1.0, 1.0],
+      [0.0, 1.0], [1.0, 0.0]];
rust> let target = vec![0, 1, 1, 0];
rust> let task = ClassificationTask::new("xor", features, target)
+      .unwrap();
rust>
rust> let mut xgb = XGBoost::new()
+      .with_n_estimators(50)
+      .with_max_depth(3);
rust> let model = xgb.train_classif(&task).unwrap();
rust> let pred = model.predict(task.features()).unwrap();
\end{CodeInput}
\end{CodeChunk}

\subsection{Implemented algorithms}

Table~\ref{tab:learners} summarizes the 27 supervised learners. Beyond standard
algorithms, \pkg{smelt-ml} includes Geographical-XGBoost
\citep{Grekousis:2025}, Oblique Forest (SPORF) \citep{Tomita+Browne+Shen:2020},
Explainable Boosting Machine (EBM), streaming Hoeffding Trees, and Dynamic
Ensemble Selection (KNORA-E).

\begin{table}[t!]
\centering
\begin{tabular}{llcc}
\toprule
Category & Algorithm & Classif. & Regress. \\
\midrule
Trees       & Decision Tree, Oblique Tree/Forest & \checkmark & \checkmark \\
Ensembles   & Random Forest, Extra Trees, Bagging, Stacking & \checkmark & \checkmark \\
            & AdaBoost, Dynamic Ensemble Selection & \checkmark & \\
Boosting    & XGBoost, LightGBM, CatBoost & \checkmark & \checkmark \\
            & Gradient Boosting, Quantile GB, EBM & \checkmark & \checkmark \\
Linear      & Linear/Logistic Regression, Ridge, Lasso, Elastic Net & \checkmark & \checkmark \\
Other       & KNN, Gaussian NB, Linear SVM, Hoeffding Tree & \checkmark & \checkmark \\
Spatial     & Geographical-XGBoost & & \checkmark \\
\bottomrule
\end{tabular}
\caption{\label{tab:learners} Supervised learners in \pkg{smelt-ml}.}
\end{table}

Additional modules provide: K-Means, DBSCAN, and Isolation Forest (clustering/anomaly);
Random Survival Forest with C-index \citep{Ishwaran+Kogalur+Blackstone:2008}
(survival analysis); Causal Forest with honest splitting
\citep{Athey+Imbens:2016, Wager+Athey:2018} (causal inference); Classifier Chains
and Regressor Chains (multi-label/output); Quantile Regression Forest; Conformal
Prediction with CQR \citep{Romano+Patterson+Candes:2019}; and SHAP-based feature
importance.


%% -- Cache-Optimized Gradient Boosting ----------------------------------------

\section{Cache-optimized gradient boosting} \label{sec:optimization}

\subsection{The histogram bottleneck}

Histogram-based gradient boosting \citep{Chen+Guestrin:2016} discretizes
continuous features into bins and accumulates gradient/hessian statistics per bin.
The key computational kernel is:

\begin{equation} \label{eq:histogram}
\text{hist}[b] \mathrel{+}= g_i, \quad b = \text{bin\_index}[i][\text{feature}]
\end{equation}

The standard row-major storage \code{bin\_index[sample][feature]} causes cache
misses when iterating over samples for a specific feature, because consecutive
memory locations correspond to different features of the same sample.

\subsection{Column-major storage with u8 packing}

Our key optimization stores bin indices in \textbf{column-major} order:
\code{bin\_index[feature][sample]}. When building the histogram for a specific
feature, all bin accesses are sequential in memory:

\begin{CodeChunk}
\begin{CodeInput}
rust> // Column-major: sequential access (cache-friendly)
rust> for &idx in node_indices {
+      let bin = bins.get_bin(feature, idx);
+      bin_g[bin] += gradients[idx];
+      bin_h[bin] += hessians[idx];
+    }
\end{CodeInput}
\end{CodeChunk}

Combined with \textbf{u8 packing} (254 real bins + 1 NaN sentinel fit in 1 byte
vs.\ 2 bytes for \code{u16}), this halves the memory bandwidth requirement for bin
access. The same \code{HistBins} structure is shared across all three gradient
boosting implementations, avoiding code duplication.

\subsection{Additional optimizations}

For \pkg{XGBoost}: automatic switch to exact greedy split finding when
$n \leq n_{\text{bins}}$, producing optimal splits for small datasets. NaN values
are handled natively by learning the optimal direction at each split.

For \pkg{CatBoost}: oblivious tree construction with bins built \textbf{once}
before the boosting loop (rather than per iteration), prefix-sum histogram
scanning, and parallel feature evaluation via \pkg{rayon}.

For \pkg{LightGBM}: Gradient-based One-Side Sampling (GOSS) with the same
column-major histogram infrastructure.


%% -- Performance Evaluation ---------------------------------------------------

\section{Performance evaluation} \label{sec:performance}

\subsection{Experimental setup}

Benchmarks use synthetic datasets from \pkg{scikit-learn}'s
\fct{make\_classification} and \fct{make\_regression} (20 features, 10 informative,
$n$ from 100 to 10{,}000). Configuration: 100 trees, max depth 6, learning rate
0.3 (\pkg{XGBoost}/\pkg{CatBoost}) or 0.1 (\pkg{LightGBM}). Single-threaded
execution. \proglang{Rust} compiled with \code{target-cpu=native} (release mode).
Official libraries via \proglang{Python} API with \code{n\_jobs=1}.

\subsection{Training time}

Table~\ref{tab:classif_perf} shows classification training times.

\begin{table}[t!]
\centering
\begin{tabular}{r rrr rrr rrr}
\toprule
 & \multicolumn{3}{c}{\pkg{XGBoost}} & \multicolumn{3}{c}{\pkg{LightGBM}} & \multicolumn{3}{c}{\pkg{CatBoost}} \\
\cmidrule(lr){2-4} \cmidrule(lr){5-7} \cmidrule(lr){8-10}
$n$ & C++ & Rust & $\times$ & C++ & Rust & $\times$ & C++ & Rust & $\times$ \\
\midrule
500   & 72  & 28  & 2.6 & 49  & 53  & 0.9 & 253 & 60  & 4.2 \\
1{,}000  & 90  & 48  & 1.9 & 104 & 88  & 1.2 & 254 & 98  & 2.6 \\
5{,}000  & 191 & 139 & 1.4 & 174 & 186 & 0.9 & 324 & 209 & 1.6 \\
10{,}000 & 274 & 253 & 1.1 & 209 & 243 & 0.9 & 445 & 350 & 1.3 \\
\bottomrule
\end{tabular}
\caption{\label{tab:classif_perf} Classification training time (ms).
Speedup $>1$ means \pkg{smelt-ml} is faster. All values single-threaded.}
\end{table}

\pkg{smelt-ml}'s \pkg{XGBoost} is consistently faster (1.1--2.6$\times$),
\pkg{CatBoost} is substantially faster (1.3--4.2$\times$), and \pkg{LightGBM}
is competitive (0.9--1.2$\times$). Similar patterns hold for regression tasks.

\subsection{Accuracy validation}

Table~\ref{tab:accuracy} validates correctness against \pkg{scikit-learn} on
standard datasets.

\begin{table}[t!]
\centering
\begin{tabular}{llccc}
\toprule
Dataset & Learner & \pkg{smelt-ml} & \pkg{scikit-learn} & $\Delta$ \\
\midrule
Iris        & Decision Tree & 0.953 & 0.953 & 0.000 \\
            & Random Forest & 0.953 & 0.967 & $-$0.014 \\
            & XGBoost       & 0.953 & 0.953 & 0.000 \\
Wine        & Random Forest & 0.983 & 0.972 & $+$0.011 \\
            & Logistic Reg. & 0.989 & 0.961 & $+$0.028 \\
            & XGBoost       & 0.971 & 0.944 & $+$0.027 \\
Breast Ca.  & Random Forest & 0.958 & 0.956 & $+$0.002 \\
            & XGBoost       & 0.970 & 0.974 & $-$0.004 \\
\bottomrule
\end{tabular}
\caption{\label{tab:accuracy} 5-fold CV accuracy on real datasets.}
\end{table}

All comparisons are within $\pm$3\% of \pkg{scikit-learn}.


%% -- Unique Features ----------------------------------------------------------

\section{Unique features} \label{sec:unique}

\subsection{Geographical-XGBoost}

Implementation of \citet{Grekousis:2025}: bi-square spatial kernel weights,
local \pkg{XGBoost} models per spatial unit, ensemble of global and local
predictions with adaptive $\alpha$ weight. First implementation outside the
original \proglang{Python} library.

\subsection{Causal Forest}

Honest splitting following \citet{Athey+Imbens:2016} and \citet{Wager+Athey:2018}.
Estimates conditional average treatment effects (CATE) $\tau(x) = E[Y(1) - Y(0) | X = x]$
with 95\% confidence intervals. First implementation in a \proglang{Rust} ML framework.

\subsection{Conformal prediction}

Split conformal prediction for regression and classification with conformalized
quantile regression \citep[CQR;][]{Romano+Patterson+Candes:2019}. Provides
distribution-free prediction intervals with guaranteed coverage.


%% -- Summary ------------------------------------------------------------------

\section{Summary and discussion} \label{sec:summary}

We have demonstrated that a pure-\proglang{Rust} implementation of gradient
boosted trees can match or exceed the performance of highly optimized
\proglang{C++} libraries through careful attention to memory layout and cache
utilization. The performance advantage stems from \proglang{Rust}'s ownership
model, which guarantees no aliasing (\code{\&mut} is unique) and enables LLVM to
apply more aggressive auto-vectorization than is possible with \proglang{C++},
where pointer aliasing must be assumed unless explicitly annotated.

\pkg{smelt-ml} provides a comprehensive ML toolkit for \proglang{Rust} with
unique capabilities in spatial ML, causal inference, and uncertainty
quantification. The framework contains zero \code{unsafe} blocks across
13{,}377 lines of code, demonstrating that memory safety and high performance
are not mutually exclusive.

Limitations include the performance gap on datasets beyond 10{,}000 samples
(where \proglang{C++} libraries employ SIMD intrinsics and GPU backends),
the absence of GPU support, and the simplified ordered boosting in \pkg{CatBoost}.

The software is available at \url{https://crates.io/crates/smelt-ml} under the
MIT license, with source code at \url{https://github.com/franciscoparrao/smelt}.


%% -- Computational details ----------------------------------------------------

\section*{Computational details}

The results in this paper were obtained using \proglang{Rust}~1.89 with the
\pkg{ndarray}~0.16, \pkg{rayon}~1.10, and \pkg{serde}~1.0 crates. Official
\proglang{C++} libraries were accessed via \proglang{Python}~3.12 with
\pkg{xgboost}~3.1, \pkg{lightgbm}~4.6, and \pkg{catboost}~1.2.
\pkg{scikit-learn}~1.8 was used for accuracy validation.
\proglang{Rust} code was compiled with \code{RUSTFLAGS="-C target-cpu=native"}
in release mode.


\section*{Acknowledgments}

[To be completed]


%% -- Bibliography -------------------------------------------------------------

\bibliography{refs}


\end{document}
